\documentclass[12pt, a4paper]{article}

\usepackage[ngerman]{babel}
\usepackage{inputenc}
\usepackage[T1]{fontenc}

\usepackage{amsmath, amssymb, physics, siunitx, amsthm}

\usepackage[left=4cm, right=3cm, head=18pt, top=2cm, bottom=2cm,includefoot=true]{geometry}

\usepackage{hyperref}
\usepackage{scrlayer-scrpage}
\usepackage{setspace}
\usepackage{changepage}

\title{Mitschriften}
\author{Maurice-León Schwinger}
\date{}

\begin{document}
	\thispagestyle{empty}
	\maketitle
	\tableofcontents
	\newpage
	\setcounter{page}{1}
	\section{Lineare Algebra}
	\subsection{Grundlagen der Mathematik}

	Aussagenlogik
	Eine Definition ist eine endliche Zeichenfolge, der genau ein Wahrheitswert, wahr oder falsch, zugeordnet ist.\\
	Zwei Aussagen $p$ (Hypothese) und $q$ (Konklusion) können wie folgt zu neuen Aussagen verknüpft werden:
	\begin{itemize}
		\item Konjunktion: $p \land q$ (min)
		\item Disjunktion: $p \lor q$ (max)
		\item Negation: $\neg p$
		\item Implikation: $p \implies q$
		\item Äquivalenz: $p \iff q$
	\end{itemize}


	Prädikatenlogik
	Ein Prädikat ordnet jedem Element $x$ aus einem Universum $X$ eine Aussage $P(x)$ zu.
	funktioniert wie Aussage

	Quantoren:
	\begin{itemize}
	\item Allquantor $\forall x \in X$
	\item Existenzquantor $\exists x \in X$
	\end{itemize}
	
	Nützliche Rechenregeln

	Mengentheoretische Grundlagen

	Mengen bestehen aus Elementen ($x \in  M$)

	Naive Definition nach Cantor: Menge Zusammenschluss von bestimmten, wohlunterscheidbaren Objekten zu einem Ganzen. 	

	Leere Menge $\emptyset$

	Beschreibung von Mengen durch Eigenschaften $A:=\{x \in B| P(x)\}$

	Mengentheoretische Bezeichnungen und Operationen

	Seien A und B Mengen.

	\begin{itemize}
	\item Enthaltenseinsbeziehung (Inklusion):\\ A heißt Teilmenge von B ($A \subseteq B$), falls jedes Element von A auch Element von B ist. 
	\item Gleichheit ($A = B$):\\
		 $A\subseteq B \land B \subseteq A$\\
		 echte Teilmenge $A \neq B \land A \subseteq B:A \subset B$ 
	\item Schnitt $A \cap B$\\
	A und B disjunkt: $A \cap B = \emptyset$
	
	\item Vereinigung $A \cup B$
	\item Differenz	$A \setminus B$
	\item Symmetrische Differenz: $A \bigtriangleup B := (A \cup B)\setminus (A\cap B)$
	\item Potenzmenge von A: $\mathcal{P}(A) := \{C|C\subset A\}$
	\item Kartesisches Produkt: $A \times B := {(a,b)| a\in A, b\in B}$ (Tupel)

	\end{itemize}

	Kardinalität: $|A|$ Anzahl Elemente einer (endlichen) Menge

	Russelsches Paradoxon

	
	\newpage
	\section{Analysis}
	
	\newpage
	\section{Grundlagen der Informatik}

	\newpage
	\section{Technische Mechanik}

	\newpage
	\section{Eingebettete Systeme}

	\newpage


	
\end{document}
		
